\section{Conclusion}\label{sec:conclusion}
We studied the finite-key behaviour of BB84 for error statistics generated by a partial intercept-resend attack in a noisy channel. The QBER obeys $Q=\ech+(1-2\ech)p/4$, which a bit-level simulation reproduces. Combined with a composable key-length bound, this yields explicit values for the key rate, the minimum block length, and the maximal tolerable attack fraction. With $\epssec=10^{-10}$ and $\ech=0.01$, a positive key at $p=0.1$ requires $N\gtrsim1.2\times10^{4}$, and the asymptotic tolerance is approached only for $N\gtrsim10^{8}$. The general-attack bound is conservative against an IR-only adversary by a large margin ($p^\ast\approx0.36$ versus $0.58$), a reminder that security statements must be made against the strongest adversary, while performance targets may be set against realistic ones. Future work will extend the analysis to decoy-state implementations with loss and dark counts, to basis-dependent noise, and to collective attacks tailored to partial information.
